\chapter{CREST conformer ensemble (menu 43)}
\label{ch:menu43}
\menulabel{43}

\section{Purpose}

CREST samples the low-energy conformer--rotamer space on the GFN2 surface.
This menu reads the final ensemble back as a sorted table with Boltzmann
weights -- the candidate-selection step that decides which conformers are
worth a target-level re-optimisation.

\section{What you need}

The run directory of a finished CREST search (the three files
\code{crest\_conformers.xyz}, \code{crest.energies} and
\code{crest\_best.xyz}); the ensemble file path alone is accepted too.

\section{Walkthrough}

The example reads the n-butane quick search (the anti/gauche pair): the
frame energies reproduce the relative-energy table to within the
three-decimal print precision, and the Boltzmann weights at 298.15~K
split 0.7322/0.2678.

\lstinputlisting[style=fbkcode, firstline=52, lastline=64,
  title={The menu-43 example (the two-conformer ensemble with its
  cross-check and weights)}]
  {examples/expected/m43_crest.txt}

\section{What you get}

The report \code{crest\_ensemble.fbk.md} in the run directory: the table
(index / absolute energy in Eh / relative energy in kcal~mol$^{-1}$ /
Boltzmann weight at 298.15~K), the best-conformer index, and the
cross-check line between the frame energies and \code{crest.energies}.

\section{Conventions, cross-checks and boundaries (measured)}

\begin{readbox}
\begin{itemize}
  \item \textbf{Weights} use $w_i \propto \exp(-\Delta E_i / RT)$ with
    $R = 1.9872042586\times10^{-3}$~kcal~mol$^{-1}$~K$^{-1}$ (CODATA),
    normalised over the listed ensemble.
  \item \textbf{Indices are post-deduplication positions} (CREGEN's RMSD
    and energy criteria), not sampling order.
  \item \textbf{The weights are GFN2-level selectors}, not final
    stabilities: re-optimise the leading conformers at the target level
    for quantitative comparison.
  \item \textbf{An unpreoptimised input may abort} on a topology change
    during CREST's initial optimisation (the capture offers options
    A/B/C; measured on the first probe) -- such a run leaves no ensemble
    and this menu reports the missing file with the next step.
\end{itemize}
\end{readbox}
